\bookchapter{الفصل 13: توصيات تطبيقات الأدوات}

\chapterintro{يقدم هذا الفصل أدوات مختارة للإنتاجية والصحة والسفر، مع تكملة قدرات النظام المفيدة والأجهزة المتعددة: أدوات تنظيف مثل Qingzhuo ومنع الإزعاج، ومفاتيح الحظر المدمجة في الهاتف، والطقس وتحذيرات الزلازل؛ التحكم عن بُعد في الكمبيوتر عبر Sunlogin/ToDesk، ومكتب الذكاء الاصطناعي مثل Trae، وسطح المكتب الاصطناعي، والنقل السريع على الشبكة المحلية مثل LocalSend، وحافظة عبر الأجهزة، ورموز Aegis/Microsoft Authenticator، وجدولة TickTick والمسح الورقي.}

\booksection{أولاً: أدوات الإنتاجية: إدارة وقتك وانتباهك}

أدوات الإنتاجية هي «الهيكل» للحياة الرقمية، وتساعدك على تنظيم الأفكار والمهام والمواد المتفرقة. إليك التوصيات حسب السيناريوهات.

\booksubsection{قائمة المهام (إدارة المهام)}

\begin{itemize}
  \item TickTick (دي دا تشينغدان) — لمن؟ لمن يحتاج إلى تسجيل مهام العمل والحياة والدراسة ومواعيدها جميعها في مكان واحد؛ وهو أحد التطبيقات الموصى بها بشدة في هذا القسم. أهم المزايا: مزامنة عبر الهاتف/الحاسوب/الجهاز اللوحي/الويب؛ ودمج القوائم والتقويم ومؤقت البومودورو في تطبيق واحد؛ وإدخال بلغة طبيعية مثل «اجتماع غدًا الساعة 3 عصرًا» و«تسليم التقرير الأسبوعي كل اثنين» مع التعرف تلقائيًا على الوقت وقواعد التكرار؛ ويدعم المهام الفرعية والوسوم والأولويات والمرفقات وقوائم العمل المشترك. مناسب لإدارة متكاملة تجمع «تنظيم المواعيد + تسجيل المهام + المراجعة اليومية».
  \item كيفية استخدام TickTick بكفاءة أكبر: صندوق التجميع: ضع كل شيء أولًا في «اليوم/صندوق التجميع»، ثم صنّفه مساءً إلى قوائم العمل والعائلة والدراسة.
  \item عرض التقويم: اسحب الاجتماعات ومواعيد التسليم إلى التقويم لرؤية التعارضات، وهو أوضح من قائمة المهام فقط.
  \item المهام المتكررة: اضبط تناول الدفع والفواتير والتمارين كمهام ذات قواعد، دون تعبئة يدوية في كل مرة.
  \item مؤقت بومودورو: ابدأ التركيز مباشرة من المهمة، واحصِ «ما أنجزته فعلًا هذا الأسبوع».
  \item التكامل مع ميزات الهاتف: قل للمساعد الصوتي «ذكّرني...» لإضافتها إلى تذكيرات النظام؛ وفعّل التقويم أيضًا للمواعيد المهمة كضمانة إضافية.
\end{itemize}

Todoist الفئة المستهدفة: المستخدمون الذين يفضلون البساطة والعمل عبر المنصات. المزايا الأساسية: واجهة نظيفة، وبنية مشاريع واضحة (مهمة مهمة فرعية مهمة فرعية ثانوية)، ونقاط karma لتحفيز بناء العادات.

Microsoft To Do الفئة المستهدفة: موظفو المكاتب ومستخدمو Outlook. المزايا الأساسية: مجاني بالكامل وبدون إعلانات، وميزة «يومي» تُفرَّغ يوميًا لتخفيف العبء النفسي، وتكامل عميق مع منظومة Microsoft.

ملاحظات/تذكيرات النظام الفئة المستهدفة: المستخدمون الخفيفون الذين لا يريدون تثبيت تطبيقات طرف ثالث. المزايا الأساسية: تجربة أصلية وبدون تأخير، وإضافة مباشرة عبر Siri/Xiao Ai، وتكامل مع أدوات النظام.

\booksubsection{مؤقت بومودورو (تدريب التركيز)}

\begin{itemize}
  \item Forest الفئة المستهدفة: من لا يستطيعون التحكم في استخدام الهاتف. المزايا الأساسية: ازرع شجرة أثناء التركيز، وإذا غادرت التطبيق مبكرًا تذبل الشجرة، مما يبني تغذية راجعة إيجابية بـ«الحياة الافتراضية»، ويمكنك الزراعة مع الأصدقاء.
  \item Tide (تشاو شي) — لمن؟ لمن يحتاج إلى الضجيج الأبيض ليساعده على التركيز أو النوم. أهم المزايا: يجمع في تطبيق واحد أصوات الطبيعة (صوت المطر والأمواج والغابات) مع مؤقت البومودورو والتأمل، بواجهة بسيطة للغاية ومريحة للنفس.
  \item المبدأ: 25 دقيقة تركيز + 5 دقائق راحة كلاسيكيًا، ويمكن تعديله إلى 50/10 أو 90/20 حسب حالتك.
\end{itemize}

\booksubsection{التقويم (تصوير الوقت)}

\begin{itemize}
  \item تقويم النظام الفئة المستهدفة: جميع المستخدمين. المزايا الأساسية: إشعارات أصلية مستقرة واستهلاك منخفض جدًا للبطارية، ويمكن الإضافة عبر Siri/Xiao Ai.
  \item Google Calendar الفئة المستهدفة: مستخدمو الأجهزة المتعددة ومن يحتاجون إلى مشاركة الجدول مع الآخرين. المزايا الأساسية: ميزة «الأهداف» تساعدك تلقائيًا على جدولة التمارين/الدراسة، والتعرف الذكي على الأحداث من رسائل Gmail.
\end{itemize}

\booksubsection{تطبيقات الملاحظات (ترسيخ المعرفة)}

\begin{itemize}
  \item Notion الفئة المستهدفة: من يحتاجون إلى إدارة موحدة للملاحظات والمهام وقواعد البيانات والويكي؛ الطلاب والمستقلون والتعاون الجماعي. المزايا الأساسية: قاعدة معرفة شاملة، تحرير معياري بالكتل، يمكنك بناء صفحة شخصية وقوائم قراءة ولوحات مشاريع، مع منظومة قوالب غنية.
  \item Obsidian الفئة المستهدفة: مديرو المعرفة المتعمقون والباحثون والكتاب. المزايا الأساسية: تخزين Markdown محلي، وروابط ثنائية الاتجاه + خريطة العلاقات، لتكوين «دماغ ثانٍ» من ملاحظاتك، والبيانات ملكك بالكامل.
  \item flomo الفئة المستهدفة: من يريدون تسجيل الأفكار المتفرقة دون كتابة مقالات طويلة. المزايا الأساسية: ملاحظات بأسلوب البطاقات، تسجيل الأفكار كما لو كنت تنشر Moments، مع دعم تصنيف \# الوسوم والرجوع حسب الخط الزمني.
  \item Apple Notes الفئة المستهدفة: مستخدمو iPhone/Mac واحتياجات التسجيل الخفيفة. المزايا الأساسية: إقلاع سريع، ودعم مسح المستندات/الرسم/القوائم/مزامنة iCloud، ومجاني وتكامل عميق مع النظام.
  \item Evernote الفئة المستهدفة: المستخدمون المتقدمون الذين يحتاجون إلى قص صفحات الويب والأرشفة بكثرة. المزايا الأساسية: إضافة قص صفحات الويب قوية، وتعرف OCR على نصوص الصور، وعلامة تجارية عريقة ومستقرة، لكن النسخة المدفوعة تحتاج اشتراكًا.
\end{itemize}

\booksubsection{أدوات PDF}

\begin{itemize}
  \item PDF Expert الفئة المستهدفة: مستخدمو iPad/iPhone، ورجال الأعمال الذين يحتاجون إلى تعليقات على العقود، والطلاب. المزايا الأساسية: يُعد أقوى محرر PDF على iOS، تعليقات سلسة، ودعم التوقيع والدمج والتقسيم وملء النماذج.
  \item WPS الفئة المستهدفة: موظفو المكاتب الذين يحتاجون إلى معالجة Word/Excel/PPT/PDF معًا. المزايا الأساسية: حزمة مكتبية متكاملة، توافق جيد مع أقراص التخزين المحلية وملفات WeChat، والنسخة المجانية كافية للاستخدام اليومي.
  \item Adobe Acrobat الفئة المستهدفة: من يتطلبون معايير صارمة لتنسيق PDF (القانون والتصميم والنشر). المزايا الأساسية: المعيار الرسمي، وتحويل PDF/A دون أخطاء، ومزامنة عبر الأجهزة وتوقيع إلكتروني احترافي.
\end{itemize}

\booksubsection{أدوات المسح (رقمنة المستندات الورقية)}

العقود الورقية والفواتير والواجبات وملاحظات السبورة ونسخ الهوية... يمكن لكاميرا الهاتف أن تحل محل معظم الماسحات. العملية الأساسية: التصوير بوضوح القص والتصحيح التلقائي تصدير PDF/صورة OCR اختياري.

\begin{itemize}
  \item CamScanner الفئة المستهدفة: رجال الأعمال والطلاب الذين يحتاجون إلى مسح العقود والإيصالات والوثائق. المزايا الأساسية: التعرف التلقائي على الحواف، والتصحيح المنظوري، وتعرف OCR على النصوص، مع إنشاء PDF متعدد الصفحات ومزامنة سحابية.
  \item Microsoft Lens الفئة المستهدفة: الطلاب ومن يحتاجون إلى مسح السبورة/صفحات الكتب إلى Word. المزايا الأساسية: مجاني وبدون إعلانات، وتصدير بنقرة واحدة إلى OneNote/Word، ودقة عالية في التعرف على النصوص.
  \item مسح Apple Notes الفئة المستهدفة: مستخدمو iPhone واحتياجات المسح الخفيفة. المزايا الأساسية: لا حاجة لتنزيل تطبيق، انقر على أيقونة الكاميرا داخل الملاحظات للمسح، مع إزالة الظلال تلقائيًا والتوقيع والتعليق، ومجاني وسهل.
  \item Baimiao / CamScanner Web / كاميرا العلامة التجارية «مسح المستندات»: معظم كاميرات Huawei وXiaomi وOPPO وvivo تتضمن وضع «المستند/المسح»؛ ويمكن استخدام Google Camera أو «امسح» في Android. مناسب لتصوير صفحة عقد عند الضرورة.
  \item حيل التصوير: إضاءة متساوية، والهاتف موازٍ للورقة، ومحاذاة الزوايا الأربع؛ صوّر المستندات متعددة الصفحات بالترتيب ثم ادمجها في PDF؛ انتبه للعلامة المائية ومكان حفظ الوثائق الحساسة، ويمكن تنظيفها نهائيًا من «المحذوفة حديثًا» بعد الاستخدام.
\end{itemize}

\booksection{ثانيًا: الصحة والرياضة}

تطبيقات الصحة هي «المساعد الشخصي» الأكثر إهمالًا للهاتف — فهي تتابع نومك ونبض قلبك ونشاطك باستمرار، وتعكس حالتك طويلة المدى أكثر من زيارة المستشفى العرضية.

\booksubsection{الرياضة واللياقة}

\begin{itemize}
  \item Keep الفئة المستهدفة: المتدرّبون في المنزل والمبتدئون ومن يريدون خسارة الدهون وبناء الجسم. المزايا الأساسية: دورات مجانية ضخمة من HIIT إلى اليوغا، وخطة تدريب تولدها الذكاء الاصطناعي تلقائيًا، ومجتمع نشط.
  \item Nike Training Club (NTC) الفئة المستهدفة: من يبحثون عن إرشاد مدربين محترفين وتدريب منهجي. المزايا الأساسية: منتج Nike الرسمي، جميع الدورات مجانية، وعروض الحركات من مدربين محترفين، وتصنيف واضح حسب العضلات/الأجهزة.
  \item Apple Fitness الفئة المستهدفة: مستخدمو Apple Watch ومن يريدون تحفيز أنفسهم بالحلقات المغلقة. المزايا الأساسية: ثلاث حلقات نشاط يومية (الحركة/التمرين/الوقوف)، وأهداف بصرية، وتكامل سلس مع بيانات نبض الساعة.
\end{itemize}

\booksubsection{تتبع النوم}

\begin{itemize}
  \item AutoSleep الفئة المستهدفة: مستخدمو Apple Watch ومن يريدون معرفة نسبة النوم العميق/الخفيف كل ليلة. المزايا الأساسية: تتبع تلقائي عند ارتداء الساعة أثناء النوم، دون بدء/إيقاف يدوي، مع درجة جودة النوم ونصائح للنوم.
  \item Sleep Cycle الفئة المستهدفة: من لا يرتدون الساعة أثناء النوم ويريدون المراقبة عبر الهاتف. المزايا الأساسية: تحليل التقلب/الشخير عبر الميكروفون، واليقظة الذكية في مرحلة النوم الخفيف لتستيقظ دون معاناة.
  \item Mi Fitness الفئة المستهدفة: مستخدمو أساور/ساعات Xiaomi. المزايا الأساسية: مجاني وتكامل عميق مع سوار Xiaomi، وتتبع موحد للنبض/الأكسجين/النوم، وقيمة مقابل سعر عالية.
\end{itemize}

\booksubsection{التأمل والاسترخاء}

\begin{itemize}
  \item Calm الفئة المستهدفة: الموظفون تحت الضغط ومن يصعب عليهم النوم. المزايا الأساسية: تأمل «الهدوء اليومي» + قصص النوم (بقراءة مشاهير)، وواجهة جميلة، ومحتوى إنجليزي غني.
  \item Tide (تشاو شي) — لمن؟ للمستخدمين الذين يحتاجون إلى الضجيج الأبيض للتركيز أو النوم. أهم المزايا: تطبيق صيني مميز يجمع أصوات الطبيعة ودروس التأمل ومؤقت البومودورو، مع إرشادات باللغة الصينية أقرب إلى القلب.
  \item Little Sleep الفئة المستهدفة: من يحتاجون إلى أصوات مساعدة على النوم أو تسجيل الكلام أثناء النوم/الشخير. المزايا الأساسية: ضوضاء بيضاء ضخمة قابلة للدمج (مطر + مروحة + نار)، وإيقاف مؤقت قبل النوم، وتسجيل الكلام ممتع ومفيد.
\end{itemize}

\booksubsection{صحة المرأة}

\begin{itemize}
  \item Flo الفئة المستهدفة: من يحتاجون إلى تسجيل الدورة الشهرية أو التخطيط للحمل أو الاهتمام بصحة المرأة. المزايا الأساسية: تنبؤ ذكي بدقة عالية بالدورة/يوم الإباضة، ومجتمع للأسئلة المجهولة، وتصميم أنيق.
  \item Meiyou الفئة المستهدفة: المستخدمات الصينيات ومن يفضلن المحتوى والعلوم بالصينية. المزايا الأساسية: محتوى محلي غني، مع معرفة ملائمة لكل مرحلة (الدورة/التخطيط للحمل/الحمل/تربية الأطفال).
\end{itemize}

\booksubsection{تذكير الأدوية}

\begin{itemize}
  \item Medication Reminder الفئة المستهدفة: مرضى الأمراض المزمنة الذين يحتاجون إلى تناول أدوية طويلة الأمد، والأبناء الذين يهتمون بأدوية كبار السن. المزايا الأساسية: مسح الدواء للتعرف عليه، وإنشاء قائمة أدوية تلقائيًا، وتنبيهات في الوقت المحدد، وتسجيل تناول الدواء ومشاركته مع العائلة.
\end{itemize}

\booksubsection{مركز البيانات}

\begin{itemize}
  \item Apple Health / Google Fit الفئة المستهدفة: جميع مستخدمي الهواتف الذكية. المزايا الأساسية: تجميع كل البيانات من الهاتف والسوور وتطبيقات الصحة الأخرى (الخطوات والنبض والنوم والدورة الشهرية إلخ)، وهو «مركز السجل الصحي» الموحد.
\end{itemize}

\booksection{ثالثًا: السفر والتنقل}

أثناء السفر، بعض التطبيقات العملية توفر عليك الكثير من المتاعب. إليك التوصيات حسب مراحل الرحلة.

\booksubsection{الطيران والقطارات فائقة السرعة}

\begin{itemize}
  \item UMetrip (هانغ لؤ تشنغ هونغ) — لمن؟ لرجال الأعمال الذين يسافرون كثيرًا بالطائرة. أهم المزايا: بيانات الطيران المدني الرسمية، وهي الأكثر موثوقية في متابعة حالة الرحلات وتغييرات بوابة الصعود والتنبيه المبكر بالتأخير، مع إمكانية الاطلاع على طراز الطائرة وخريطة المقاعد.
  \item 12306 الفئة المستهدفة: جميع المسافرين بالقطارات في الصين. المزايا الأساسية: القناة الرسمية الموثوقة الوحيدة لشراء التذاكر، ودعم التذاكر البديلة (عند إلغاء حجز أحدهم يُحجز تلقائيًا)، وتسهيل تذاكر الطلاب والمجموعات.
  \item VariFlight الفئة المستهدفة: من يهتمون بدقة مواعيد الرحلات ومن يحتاجون إلى الاستقبال والتوصيل. المزايا الأساسية: تنبؤ دقيق بتأخير الرحلات، ومشاركة الرحلة مع العائلة لمعرفة موقعك ووقت وصولك لحظيًا.
\end{itemize}

\booksubsection{الخرائط والملاحة}

\begin{itemize}
  \item Amap الفئة المستهدفة: مستخدمو المدن في الصين وسائقو السيارات/الأجرة. المزايا الأساسية: حالة المرور اللحظية الدقيقة، وقدرة قوية على تجنب الازدحام، وحجز الأجرة/المشاركة في مكان واحد، وخيارات نقل عام غنية.
  \item Baidu Maps الفئة المستهدفة: من يحتاجون إلى معلومات الأماكن والملاحة الداخلية. المزايا الأساسية: خرائط بانوراما/ملاحة AR، وخرائط داخلية تفصيلية للمناطق التجارية والمتاحف، وحزم أصوات غنية.
  \item Google Maps الفئة المستهدفة: المسافرون إلى الخارج. المزايا الأساسية: أوسع تغطية عالمية، ومراجعات محلية حقيقية، والصور الشارعية والخرائط دون اتصال منقذة في البلدان الأجنبية.
\end{itemize}

\booksubsection{الترجمة}

\begin{itemize}
  \item Google Translate الفئة المستهدفة: المسافرون والدارسون عبر اللغات. المزايا الأساسية: ترجمة «المحادثة» الفورية ثنائية الاتجاه، وترجمة «الكاميرا» الفورية، وهي أدوات سفر رائعة، مع دعم حزم اللغات دون اتصال.
\end{itemize}

\booksubsection{تحويل العملات}

\begin{itemize}
  \item Minimal Exchange Rate الفئة المستهدفة: المسافرون والمتسوقون عبر الحدود. المزايا الأساسية: واجهة نظيفة للغاية، ودعم أكثر من 160 عملة، وتحديث فوري للأسعار، ومقارنة عدة عملات في وقت واحد، مع تحويل سريع بالضغط المطول على الأرقام.
\end{itemize}

\booksubsection{الطقس}

\begin{itemize}
  \item تطبيق الطقس في النظام / China Weather الفئة المستهدفة: الجميع، الخيار الأول للتنبؤات والتحذيرات الرسمية. المزايا الأساسية: يتكامل تطبيق الطقس من الشركات مع القنوات الرسمية مثل «China Weather» لاستقبال تحذيرات الأرصاد، مثل الأمطار الغزيرة والأعاصير والحرارة، عبر لافتات/إشعارات، بدرجة موثوقية عالية.
  \item Caiyun Weather الفئة المستهدفة: من يحتاجون إلى تنبؤات أمطار دقيقة على مستوى الدقيقة. المزايا الأساسية: دقة عالية في «هل ستمطر خلال الساعتين المقبلتين»، وخريطة ديناميكية للأمطار مرئية بوضوح، مع تنبيهات استباقية قبل المطر.
  \item Carrot Weather الفئة المستهدفة: هواة الطقس الذين يبحثون عن المتعة وودجتات غنية. المزايا الأساسية: تعليق صوتي ساخر بالذكاء الاصطناعي، وودجتات سطح مكتب مخصصة كثيرة، وخرائط رادار احترافية، والنسخة المدفوعة أكثر شمولاً.
  \item نقاط التحذير: يجب السماح بالإشعارات لتطبيقات الطقس لاستقبال التحذيرات؛ التحذيرات أربع درجات (أزرق وأصفر وبرتقالي وأحمر)، والأحمر يعني أعلى مخاطر للطقس العاصف، ويجب تعديل السفر حسب التعليمات. لطرق التفعيل الأخرى راجع القسم السابع من هذا الفصل «تحذيرات الطقس / الكوارث الطبيعية / الزلازل».
\end{itemize}

\booksubsection{أدلة السفر}

\begin{itemize}
  \item Mafengwo الفئة المستهدفة: من يحبون إعداد أدلة مفصلة ورحلات السيارات/السفر العميق. المزايا الأساسية: محتوى تقارير المستخدمين غني، ودليل الوجهات/المسارات/تجنب الأخطاء مفصل، ويمكن تنزيل الأدلة دون اتصال.
  \item Xiaohongshu الفئة المستهدفة: المسافرون الشباب الذين يريدون تجارب حقيقية وأماكن غير شهيرة/مأكولات مميزة. المزايا الأساسية: صور ومقاطع قصيرة بصرية، ومحتوى مثل «دليل مدينة XX لمدة 3 أيام» محدث بسرعة وذو صلة زمنية، وتجربة بحث جيدة.
\end{itemize}

\booksection{رابعًا: مدخل لأدوات الأتمتة}

الأتمتة هي رافعة «لتبديل الإعداد مرة واحدة مقابل توفير طويل المدى»، ويمكنها ضغط عشرات العمليات المتكررة يوميًا إلى الصفر. هذا القسم مدخل للمفاهيم فقط، وللتفاصيل والدروس راجع الفصل 23: التفكير البرمجي ومدخل الأتمتة.

\begin{itemize}
  \item أوامر iOS «الاختصارات»: أداة الأتمتة الرسمية من Apple، دون برمجة، يمكنها تنفيذ إجراءات مركبة بنقرة واحدة مثل «عند الوصول إلى المنزل: تشغيل WiFi + خفض السطوع + نشر الطقس».
  \item أتمتة Android: ربط السيناريوهات على مستوى النظام (مثل «المهام التلقائية» في Xiaomi و«الذكاء السياقي» في Huawei) هو الأسهل؛ أما Tasker فقوي لكن منحنى التعلم حاد، ومناسب للمحترفين.
  \item الفكرة الأساسية: حدد أولًا «العمليات المتكررة يوميًا»، ثم فكر «هل يمكن استبدالها بشرط تشغيل واحد + عدة إجراءات». للأمثلة التفصيلية والتقنيات المتقدمة راجع الفصل 23.
\end{itemize}

\booksection{خامسًا: المتصفحات ومحركات البحث: بوابتك إلى المعلومات}

المتصفح هو نافذتك إلى الإنترنت، ومحرك البحث هو بوصلتك في بحر المعلومات. استخدامهما جيدًا يوفر عليك الكثير من المنعطفات.

\booksubsection{مبدأ عمل المتصفح:}

عند إدخال عنوان الموقع في المتصفح: (1) يستعلم المتصفح من خادم DNS عن عنوان IP المرتبط بالنطاق؛ (2) يرسل طلب HTTP/HTTPS إلى خادم ذلك العنوان؛ (3) يعيد الخادم كود الصفحة (HTML وCSS وJavaScript)؛ (4) يحلل المتصفح الكود ويعرض الصفحة التي تراها.

\booksubsection{توصيات المتصفحات الرئيسية:}

\begin{itemize}
  \item Microsoft Edge: مبني على محرك Chromium، سريع ومتوافق جيدًا. يتضمن مساعد Copilot AI لتلخيص الصفحات والترجمة والمحادثة. يدعم المزامنة عبر الأجهزة، ونوصي به بشدة كمتصفح رئيسي.
  \item Safari: المتصفح الافتراضي لأجهزة Apple، متكامل بعمق مع النظام، موفر للطاقة وفعّال. حماية جيدة للخصوصية، لكن الإضافات أقل.
  \item Chrome: الأكثر استخدامًا عالميًا، إضافات غنية. لكنه يستهلك ذاكرة كبيرة، ويستهلك بطارية أكثر على الهاتف.
\end{itemize}

\booksubsection{اختيار محرك البحث:}

\begin{itemize}
  \item Bing: محرك بحث Microsoft، متكامل بعمق مع Edge. ميزة الملخص الذكي تعطي الإجابة مباشرة وليس فقط الروابط، وكفاءة البحث عالية. تجربة البحث بالصور جيدة أيضًا.
  \item Yandex: أكبر محرك بحث روسي، قوي في البحث بالصور والبحث العكسي عن الصور. مفيد جدًا للعثور على صور مشابهة والتعرف على الأشياء.
  \item Baidu: الرئيسي للبحث عن المعلومات في الصين، لكن الإعلانات كثيرة وتحتاج إلى قدرة تمييز.
  \item Google: الأدق عالميًا، لكنه يحتاج إلى شبكة خاصة للاستخدام في الصين.
\end{itemize}

\booksubsection{دليل البحث الفعال عن المعلومات:}

\begin{itemize}
  \item استخدم البحث الذكي: ملخصات Bing الذكية وCopilot في Edge تجيب مباشرة، وتستخلص المعلومات المهمة من النتائج الكثيرة، وهو أكثر كفاءة من النقر على الروابط واحدًا واحدًا.
  \item التعرف على المعلومات المزيفة والإعلانات: البنود الموسومة بـ«إعلان» أو «مُموَّل» في نتائج البحث هي ترويج مدفوع وليست نتائج طبيعية. انتبه لموثوقية المصدر، وفضل المحتوى من المواقع الرسمية والمعروفة.
  \item تنزيل التطبيقات بشكل صحيح: اكتب في محرك البحث «اسم التطبيق + الموقع الرسمي»، مثل «الموقع الرسمي لـ WeChat»، وانقر أولًا على النتائج التي تحمل النطاق الرسمي. لا تثبت البرامج من مواقع التنزيل التابعة لجهات ثالثة، فهي غالبًا ترفق برامج ضارة.
  \item حيلة تمييز المواقع الرسمية: لتنزيل Windows 11 ازور microsoft.com، ولتنزيل Chrome ازور google.com. التحقق من تطابق النطاق مع العلامة التجارية هو أبسط طريقة للتمييز.
\end{itemize}

\booksubsection{معرفة نطاقات المواقع:}

النطاق هو «عنوان الموقع»، وفهمه يساعدك على الحكم على طبيعة الموقع ومصدره:

\begin{itemize}
  \item معاني اللاحقة الشائعة: .com — مؤسسة تجارية، الأكثر شيوعًا
  \item .cn — النطاق الوطني للصين
  \item .org — منظمة غير ربحية
  \item .net — مؤسسة خدمات الشبكات
  \item .edu — مؤسسة تعليمية
  \item .gov — مؤسسة حكومية
  \item .xyz/.tech/.info — نطاقات عامة جديدة، منخفضة المتطلبات، احذر منها
\end{itemize}

نطاقات الدول/المناطق: .jp — اليابان

.kr — كوريا الجنوبية

.uk — المملكة المتحدة

.de — ألمانيا

.tw — تايوان

.hk — هونغ كونغ

تنبيه أمني: مواقع .gov و.edu عادةً أكثر موثوقية لأن متطلبات التسجيل عالية. أما النطاقات الرخيصة مثل .xyz و.top فتُستخدم كثيرًا في مواقع التصيد، فانتبه عند الزيارة.

\booksubsection{حظر الإعلانات:}

\begin{itemize}
  \item AdGuard الفئة المستهدفة: جميع المستخدمين المصابين بالإعلانات المنبثقة والداخلة في التطبيقات. المزايا الأساسية: حظر إعلانات على مستوى النظام (وليس المتصفح فقط)، وتصفية لافتات الإعلانات وتتبع الشاشات الافتتاحية داخل التطبيقات، وتوفير البيانات وتسريع تحميل الصفحات، والنسخة المدفوعة تدعم جميع الأجهزة.
\end{itemize}

\booksubsection{متصفحات الخصوصية:}

\begin{itemize}
  \item Brave الفئة المستهدفة: المتصفحون اليوميون الذين يهتمون بالخصوصية ولا يريدون التتبع الإعلاني. المزايا الأساسية: مبني على Chromium، يحظر الإعلانات والمتتبعات جاهزًا، ويتضمن وضع التخفي Tor، وأسرع من Chrome وأقل استهلاكًا للذاكرة.
  \item DuckDuckGo الفئة المستهدفة: من يبحثون عن «البحث بدون ملاحقة إعلانية». المزايا الأساسية: محرك بحث بالاسم نفسه (لا يسجل سجل البحث)، ومسح جميع آثار التصفح بنقرة واحدة («زر الحرق»)، وتجربة موحدة عبر المنصات.
\end{itemize}

\booksubsection{مقارنة المتصفحات:}

\begin{center}
\begin{tabularx}{\linewidth}{|X|X|X|}
\hline
\textbf{المتصفح} & \textbf{المزايا الأساسية} & \textbf{الفئة المستهدفة} \tabularnewline
\hline
Microsoft Edge & محرك Chromium + Copilot AI مدمج، يمكنه تلخيص الصفحات/الترجمة/المحادثة & موظفو المكاتب الباحثون عن الكفاءة ومستخدمو الذكاء الاصطناعي المكثفون (موصى به كرئيسي) \tabularnewline
Safari & تكامل عميق مع منظومة Apple، موفر للطاقة وفعّال، وحماية قوية للخصوصية & مستخدمو iPhone/Mac، المهتمون بالبطارية والتكامل مع النظام \tabularnewline
Chrome & الأكثر استخدامًا عالميًا، وأغنى منظومة إضافات & المستخدمون المتقدمون الذين يحتاجون إلى إضافات كثيرة، ومحبو المزامنة عبر الأجهزة \tabularnewline
Brave & يحظر الإعلانات والمتتبعات جاهزًا، وأسرع & المهتمون بالخصوصية الذين يريدون حظر الإعلانات دون إضافات \tabularnewline
DuckDuckGo & محرك بحث بالاسم نفسه لا يسجل التاريخ، ومسح بنقرة واحدة & المستخدمون الحساسون للخصوصية الذين لا يريدون التتبع المخصص \tabularnewline
\hline
\end{tabularx}
\end{center}

\begin{tcbpractice}{جرّب إنشاء سير عمل آلي «لاسترخاء قبل النوم»}
\begin{enumerate}
  \item أنشئ شرط التشغيل: وقت محدد (مثل 22:30).
  \item أضف الإجراءات: خفض الصوت > تفعيل وضع عدم الإزعاج > تشغيل قائمة استرخاء > خفض السطوع.
  \item احفظ واختبر.
\end{enumerate}

\end{tcbpractice}

\booksection{سادسًا: أدوات النظام المفيدة: التنظيف ومنع الإزعاج والحماية}

إلى جانب تطبيقات الإنتاجية والصحة، غالبًا ما تؤثر ثلاثة أمور على التجربة اليومية: «تقليل الإزعاج، ومساحة تخزين كافية، ومنع المحتالين». يتناول هذا القسم أولًا البرامج العملية الشائعة، ثم وظائف الحظر المدمجة في الهاتف — فقد اشترى كثيرون هواتفهم منذ سنوات دون فتح هذه المفاتيح أبدًا.

\booksubsection{تنظيف النظام وتحسين التخزين}

\begin{itemize}
  \item QingZhuo (تشينغ جوه) — لمن؟ لمستخدمي Android الذين يرغبون في تنظيف خفيف للذاكرة المؤقتة والملفات المتبقية والمجلدات الفارغة دون أن يزعجهم «حزمة مدير الهاتف» الكاملة. أهم المزايا: حجم صغير وبلا إعلانات مزعجة، ويدعم الفحص العميق للبقايا التي تخلفها التطبيقات وتحديد مواقع الملفات الكبيرة وتحليل استهلاك مساحة التخزين؛ وهو أكثر تحفظًا من أدوات «التسريع بضغطة واحدة» التجارية، ومناسب للمستخدمين أصحاب المهارة اليدوية للصيانة الدورية. تنبيه: تأكد من المسارات قبل التنظيف، ولا تحدد مجلدات WeChat والألبوم والمستندات كلها دفعة واحدة دون تمييز؛ واجعل التنظيف المدمج في هواتف العلامة التجارية هو الخيار الأول للذاكرة المؤقتة على مستوى النظام، واكتفِ بالأدوات الخارجية كمكمل فقط.
  \item أداة التنظيف المدمجة في النظام الفئة المستهدفة: الجميع، خاصة كبار السن والمستخدمون الذين يخشون التعقيد. المزايا الأساسية: Android [الإعدادات] > [التخزين] > [تنظيف وتسريع]؛ iOS [الإعدادات] > [عام] > [مساحة تخزين iPhone]. النظام يفهم بنية الذاكرة المؤقتة لتطبيقاته بشكل أفضل، وخطر الحذف الخاطئ هو الأدنى، لذا يجب أن يكون الخيار الأول.
  \item إدارة التخزين المدمجة في WeChat/QQ الفئة المستهدفة: المستخدمون الذين يملكون محادثات كثيرة. المزايا الأساسية: WeChat [أنا] > [الإعدادات] > [عام] > [مساحة التخزين]، يمكن التنظيف الدقيق حسب المحادثة ونوع الملف؛ أكثر أمانًا من أدوات «تنظيف WeChat» الخارجية.
  \item المبدأ: ما يمكن تنظيفه بالنظام/التطبيق المدمج فلا تستخدم الأدوات التابعة؛ وما يمكن تحديده حسب المسار/النوع فلا تستخدم «التنظيف الشامل بنقرة واحدة». الهدف هو تحرير المساحة، وليس مسح البيانات المهمة.
\end{itemize}

\booksubsection{أدوات منع الإزعاج والحماية}

\begin{itemize}
  \item National Anti-Fraud Center الفئة المستهدفة: الجميع، خاصة كبار السن والوظائف المالية ومن يتسوقون كثيرًا أو يتلقون مكالمات غريبة. المزايا الأساسية: منتج رسمي، مع تحذيرات المكالمات وتنبيهات مخاطر الرسائل والروابط المشبوهة؛ مرتبط بقاعدة بيانات الشرطة لمكافحة الاحتيال، وهو أكثر حلقة موثقة رسميًا في «الخطوط الدفاعية الثلاثة».
  \item مركز أمان العلامة التجارية / مدير الهاتف الفئة المستهدفة: مستخدمو العلامة التجارية المعنية. المزايا الأساسية: Huawei Phone Manager وXiaomi Security Center وOPPO Phone Manager وvivo Phone Manager وغيرها، تجمع حظر الإزعاج وفحص الفيروسات وإدارة الأذونات وحماية الدفع، وعادةً مجانية وتنسق أفضل مع النظام، ولا يُنصح بتثبيت مجموعة أمان تجارية ثانية تسبب تعارض الأذونات.
  \item خدمة منع الإزعاج لدى شركة الاتصالات الفئة المستهدفة: المستخدمون الذين يقصفون بالمكالمات عالية التردد ومكالمات الذكاء الاصطناعي. المزايا الأساسية: China Mobile أرسل KTFSR إلى 10086؛ China Telecom تابع «Tianyi Anti-Harassment»؛ China Unicom اضبط في «China Unicom Mobile Manager». الحظر يتم على مستوى شبكة الاتصالات، ويمكنه صد دفعة من الأرقام المزعجة حتى دون رنين الهاتف.
  \item AdGuard الفئة المستهدفة: جميع المستخدمين المصابين بالإعلانات المنبثقة والداخلة في التطبيقات. المزايا الأساسية: حظر إعلانات على مستوى النظام (وليس المتصفح فقط)، وتصفية لافتات الإعلانات وتتبع الشاشات الافتتاحية داخل التطبيقات، وتوفير البيانات وتسريع تحميل الصفحات، والنسخة المدفوعة تدعم جميع الأجهزة.
  \item احذر تثبيت حزم «التنظيف والتسريع / حظر الإزعاج» الشاملة: بعض أدوات الإدارة التابعة تبقى في الخلفية وتطلب أذونات جهات الاتصال والرسائل باستمرار، وقد لا تكون قدرة الحظر أقوى من النظام، بل تزيد استهلاك البطارية ومخاطر الخصوصية. فضّل أولًا تفعيل الميزات المدمجة، ثم فكر في تطبيق مكافحة الاحتيال الرسمي.
\end{itemize}

\textbf{مبادئ اختيار البرامج العملية}\par

\begin{itemize}
  \item النظام أولًا: الحظر والتنظيف وإدارة الأذونات، فعّل أولًا مفاتيح الهاتف المدمجة.
  \item الإكمال الرسمي: National Anti-Fraud Center وخدمات منع الإزعاج لدى شركة الاتصالات قنوات رسمية، يمكن إضافتها.
  \item الاعتدال مع الأدوات التابعة: يكفي تثبيت أداة واحدة للوظيفة نفسها، لتجنب تعارض الأذونات والتنافس على الموارد.
  \item القليل خير من الكثير: كلما زادت الأدوات، زادت الإشعارات وسطح الأذونات، وارتفعت احتمالية الإزعاج من «الأداة نفسها».
\end{itemize}

\booksubsection{وظائف الحظر المدمجة في الهاتف: كيفية التفعيل لكل علامة}

فيما يلي المداخل حسب العلامة التجارية. قد تختلف الأسماء قليلًا حسب إصدار النظام، ويمكن البحث مباشرة في الإعدادات عن «مزعج» أو «حظر» أو «مكالمات غير معروفة». يكمل هذا القسم الفصل 8 «التعامل مع المكالمات المزعجة»: هنا نركز على مواقع المفاتيح، وهناك نركز على وسم الأرقام والاعتراض على الوسم الخاطئ.

\begin{center}
\begin{tabularx}{\linewidth}{|X|X|X|}
\hline
\textbf{المنصة} & \textbf{مسار الدخول} & \textbf{الخيارات الموصى بتفعيلها} \tabularnewline
\hline
iPhone / iOS & {}[الإعدادات] > [الهاتف] / [الرسائل] & كتم المكالمات غير المعروفة؛ جهات الاتصال المحظورة؛ «تصفية المرسلين غير المعروفين» في الرسائل \tabularnewline
Xiaomi / Redmi & تطبيق الهاتف > الإعدادات > حظر المزعجات؛ أو مركز الأمان & الحظر الذكي، وقطع المكالمات غير المعروفة بعد رنينة واحدة، وتصفية الرسائل بالكلمات المفتاحية \tabularnewline
Huawei / Honor & الهاتف > حظر المزعجات؛ مدير الهاتف > حظر المزعجات & التعرف على المتصل، والقائمة السوداء، والحظر الذكي للرسائل \tabularnewline
OPPO / realme / OnePlus & تطبيق الهاتف > حظر المزعجات؛ أو هذا الجهاز > حظر المزعجات & حظر المكالمات، والتعرف على المكالمات الاحتيالية، وقواعد حظر الرسائل \tabularnewline
vivo / iQOO & الهاتف > حظر المكالمات المزعجة؛ i Manager & حظر المزعجات، وتحذيرات المكالمات الاحتيالية، وتصفية الرسائل \tabularnewline
Samsung & الهاتف > الإعدادات > حظر الأرقام / معرفة المتصل والمكالمات المزعجة & حظر الأرقام غير المعروفة، والتعرف على المكالمات المزعجة (يختلف الدعم حسب المنطقة) \tabularnewline
\hline
\end{tabularx}
\end{center}

\begin{tcbinfo}{توصية التكامل ثلاثي الطبقات}
اليومي الموصى به: «حظر النظام + منع الإزعاج لدى شركة الاتصالات + National Anti-Fraud Center» ثلاث طبقات: النظام يتولى الوسم والقائمة السوداء، وشركة الاتصالات تحظر المكالمات عالية التردد على مستوى الشبكة، ومركز مكافحة الاحتيال يكمل التحذيرات الرسمية. لا حاجة لتفعيل أقصى صرامة في كل الطبقات، لتجنب حظر مكالمات التوصيل والطعام والأطباء — يمكن السماح لجهات الاتصال والاتصالات الأخيرة.

\end{tcbinfo}

\begin{tcbpractice}{تطبيق عملي: فعّل حظر المزعجات اليوم}
\begin{enumerate}
  \item افتح إعدادات تطبيق الهاتف، وابحث عن «حظر المزعجات / كتم المكالمات غير المعروفة»، وفعّل الحظر الذكي.
  \item افحص هاتف الوالدين بنفس الطريقة: المسار نفسه + تفعيل منع الإزعاج لدى شركة الاتصالات (China Mobile أرسل KTFSR إلى 10086).
  \item ثبّت «National Anti-Fraud Center» من متجر التطبيقات، وسجّل الدخول وفعّل تحذيرات المكالمات.
  \item اختبر رقمًا موسومًا (أو راجع سجل الحظر) للتأكد من عمل الميزة.
\end{enumerate}

\end{tcbpractice}

\booksection{سابعًا: تحذيرات الطقس والكوارث الطبيعية والزلازل}

حظر المزعجات يمنع «الاحيال الهاتفي»، وميزات التحذير تمنع «الكوارث الطبيعية والطقس القاسي». كثيرون يثبّتون تطبيق الطقس فقط لرؤية درجة الحرارة، دون تفعيل الإشعارات أو تحذيرات الزلازل — ففي الثواني القليلة قبل الأمطار الغزيرة أو الإعصار أو الزلزال، يمكن للهاتف أن يمنحك وقتًا ثمينًا للنجاة.

\booksubsection{تحذيرات الطقس: كيفية استقبال تحذيرات الأمطار والأعاصير والحرارة}

\begin{itemize}
  \item القنوات الرسمية أولًا: تحذيرات China Meteorological Administration / Central Meteorological Observatory / مكاتب الأرصاد الإقليمية هي الأكثر موثوقية. تابع حساب «China Weather» على WeChat وWeibo، أو ثبّت تطبيقات رسمية مثل «China Weather».
  \item تطبيقات الطقس المدمجة: Huawei Weather وXiaomi Weather وOPPO Weather وvivo Weather وتطبيق «الطقس» من Apple وغيرها، معظمها يرسل تحذيرات المحطة المحلية؛ ادخل إعدادات تطبيق الطقس وتأكد من تفعيل إشعارات «تحذيرات الأرصاد / الطقس العاصف».
  \item تطبيقات أمطار دقيقة مثل Caiyun Weather: مناسبة لتقييم «هل آخذ المظلة خلال نصف ساعة؟» قبل التنقل أو استقبال الأطفال؛ مع بقاء التحذيرات الرسمية هي المرجع.
  \item معاني ألوان التحذير (عام): أزرق: انتبه واحذر، سافر كالمعتاد وتابع التحديثات.
  \item أصفر: قوِّ الاحتياطات، وقلل الأنشطة الخارجية غير الضرورية.
  \item برتقالي: خطير، ابقَ في الداخل قدر الإمكان، وابتعد عن المناطق الخطرة.
  \item أحمر: خطير للغاية، اتبع تعليمات إيقاف الدراسة والعمل والنجاة، ولا تجبر على عبور الطرق الغارقة أو السيول أو الرياح القوية.
\end{itemize}

أنواع التحذيرات الشائعة: أمطار غزيرة، ثلوج غزيرة، إعاصير، حرارة عالية، موجة برد، رياح قوية، صواعق، برد، تجمد الطرق، ضباب كثيف، جفاف، مخاطر الكوارث الجيولوجية.

\booksubsection{تحذيرات الكوارث الطبيعية: ليست «المطر» فقط}

\begin{itemize}
  \item الجهات المصدرة: مكتب الأرصاد مسؤول عن التحذيرات الجوية؛ أما الكوارث الجيولوجية والسيول فتُصدرها غالبًا إدارات الموارد الطبيعية والمياه والطوارئ بشكل مشترك؛ وتُصدر الحكومات المحلية عبر الإذاعات الطارئة والرسائل و12379 وغيرها.
  \item كيفية الاستقبال عبر قنوات متعددة: إشعارات نظام الهاتف + تحذيرات تطبيق الطقس.
  \item تابع في WeChat حسابات «إدارة الطوارئ» و«مكتب الأرصاد» المحلية.
  \item قنوات 12379 الوطنية لنشر التحذيرات (رسائل/خدمات، حسب توفر الخدمة في المنطقة).
  \item التلفزيون والإذاعة الطارئة ومكبرات الصوت في القرى/الأحياء — في المناطق الجبلية وموسم الفيضانات لا تعتمد على الهاتف وحده.
\end{itemize}

تذكيرات خارجية: في الأيام العاصفة لا تقترب من لوحات الإعلانات ولا تمشي تحت الأشجار؛ وفي الأمطار الغزيرة لا تدخل مواقف السيارات والأنفاق والجسور المنخفضة؛ وفي السياحة الجبلية عند تحذير «مخاطر الكوارث الجيولوجية» غادر المنحدرات والأخاديد فورًا.

الترابط مع دليل الإسعافات: عند الانقطاع بعد الكارثة أو نفاد بطارية الهاتف أو الحاجة إلى تحديد الموقع للنجدة، راجع الفصل 32 «سيناريوهات الزلازل والكوارث»، والفصل 7 «الاتصالات عبر الأقمار الصناعية والاتصال دون شبكة».

\booksubsection{تحذيرات الزلازل: ليست «التنبؤ بالزلازل»}

تحذير الزلازل (EEW) هو إرسال إنذار إلى المناطق المحتملة التأثر قبل وصول موجات الزلازل (خاصة الموجات العرضية الأكثر تدميرًا) بثوانٍ إلى عشرات الثواني، اعتمادًا على أن الموجات الكهرومغناطيسية أسرع من موجات الزلازل، وأن المحطات تلتقط الموجات الطولية أسرع. لا يمكنه التنبؤ بحدوث الزلازل، لكنه يمنحك نافذة «للنجاة القريبة وإطفاء مصادر الحريق والخروج من المصاعد». كلما اقتربت من مركز الزلزال، قصر وقت التحذير، وقد لا يكفي.

\begin{itemize}
  \item الحلول الرسمية والصناعية: في الصين يشمل ذلك قدرات China Earthquake Early Warning Network، وخدمات التحذير التي تقدمها معاهد مثل Chengdu Institute of Disaster Mitigation؛ وبعض الهواتف تأتي مزودة مسبقًا أو تستقبل التحذير عبر الإنترنت.
  \item تطبيقات الطرف الثالث: ابحث في متجر التطبيقات عن تطبيقات «تحذير الزلازل» الرسمية/المؤسسية (راجع المطور قبل التثبيت، وتجنب تطبيقات «التحذير المزيف» التي تطلب أذونات غير ضرورية).
  \item ماذا تفعل عند استلام الإنذار: في الداخل احتمِ قريبًا (تحت الطاولة، في زاوية الجدار الحامل، بجانب أثاث متين)، واحمِ رأسك ورقبتك؛ ابتعد عن النوافذ والجدران الخارجية والأثاث العالي؛ لا تندفع عشوائيًا على الدرج أو تقفز؛ إذا كنت في المصعد اخرج في أقرب طابق؛ في الخارج اذهب إلى مكان مفتوح وابتعد عن الأسلاك واللوحات والمباني المتهدمة. بعد الزلزال افحص الغاز والكهرباء وانتبه للهزات الارتدادية.
  \item لا تكتم ولا توقف: يجب السماح لإنذار الزلازل «بكسر» وضع الصامت/عدم الإزعاج؛ ضع الهاتف في متناول اليد قبل النوم، ولا تجعل صوت الجرس صفرًا.
\end{itemize}

\booksubsection{تحذيرات الزلازل المدمجة في الهاتف: كيفية التفعيل لكل علامة}

قد تختلف المداخل التالية قليلًا حسب إصدار النظام، ويمكن البحث في الإعدادات عن «زلزال» أو «تحذير» أو «تحذير أمني». يدعم iOS تحذيرات الزلازل في الصين بشكل محدود، وعادةً يحتاج إلى تطبيق تحذير موثوق مع السماح بتحديد الموقع والإشعارات.

\begin{center}
\begin{tabularx}{\linewidth}{|X|X|X|}
\hline
\textbf{المنصة} & \textbf{المداخل الشائعة} & \textbf{الوصف} \tabularnewline
\hline
Xiaomi / Redmi & مدير الهاتف / مركز الأمان > تحذير الزلازل؛ أو ابحث في الإعدادات عن «زلزال» & بعض الإصدارات تتضمن خدمة تحذير مدمجة، وتحتاج إلى اتصال بالإنترنت وتفعيل تحديد الموقع \tabularnewline
Huawei / Honor & الإعدادات > الأمان > إشعارات التحذير الطارئ / تحذير الزلازل & قد يُسمى «إشعارات التحذير الطارئ»، ويشمل الزلازل وغيرها \tabularnewline
OPPO / realme / OnePlus & الإعدادات > الأمان (اتصال SOS الطارئ) > تحذير الزلازل & بعض الطرازات تعتمد على تغطية الخدمة في المنطقة \tabularnewline
vivo / iQOO & تطبيق الطقس أو الإعدادات > الأمان/المزيد > تحذير الزلازل & استند إلى القوائم الفعلية لإصدار النظام \tabularnewline
iPhone / iOS & عادةً لا يوجد مفتاح مدمج كامل في إعدادات النظام؛ ثبّت تطبيق تحذير زلازل موثوقًا واسمح بالإشعارات & يحتاج التطبيق إلى العمل في الخلفية وصلاحيات تحديد الموقع؛ ولا تمنعه إشعارات وضع عدم الإزعاج \tabularnewline
Samsung وغيرها & الإعدادات > الأمان والطوارئ / إنذارات الطوارئ اللاسلكية & تختلف تغطية «إنذارات الطوارئ» حسب المنطقة وشركة الاتصالات \tabularnewline
\hline
\end{tabularx}
\end{center}

\begin{tcbwarn}{مفاهيم خاطئة شائعة عن تحذيرات الزلازل}
\begin{itemize}
  \item التحذير $\neq$ التنبؤ: لا يخبرك «هل سيحدث زلزال غدًا»، بل ينذر فقط بعد وقوع الزلزال وقبل وصول الموجات.
  \item ليس كل مكان يحصل على ثوانٍ كثيرة: بالقرب من مركز الزلزال قد تكون بضعة ثوانٍ فقط أو لا تكفي، والتحذير يفيد المناطق المحيطة بشكل رئيسي.
  \item إيقاف الإشعارات يعني عدم التثبيت: حتى أفضل نظام تحذير، إذا كتمت الصوت أو قتلت التطبيق في الخلفية أو منعت تحديد الموقع، فقد لا يصلك.
  \item تطبيقات مزيفة: احذر برامج «التنبؤ بالزلازل» التي تطلب أذونات جهات الاتصال والرسائل، وفضل الحلول المدمجة أو الرسمية.
\end{itemize}

\end{tcbwarn}

\begin{tcbinfo}{كيفية دمج قدرات التحذير}
التوصية اليومية: «إشعارات تحذير الطقس في النظام + القنوات الرسمية للأرصاد/الطوارئ + تحذير الزلازل (النظام أو تطبيق موثوق)». كما في منع الإزعاج، لا حاجة لتثبيت مجموعة برامج متشابهة؛ المهم هو تفعيل الإشعارات والسماح بتحديد الموقع وإمكانية كسر وضع عدم الإزعاج. بعد إعداد حظر المزعجات لهاتف الوالدين، افحص أيضًا تحذيرات الطقس والزلازل.

\end{tcbinfo}

\begin{tcbpractice}{تطبيق عملي: افحص إعدادات التحذير في 10 دقائق}
\begin{enumerate}
  \item افتح تطبيق الطقس في النظام، ثم الإعدادات، وفعّل إشعارات «تحذيرات الأرصاد / الطقس العاصف».
  \item ابحث في إعدادات النظام عن «زلزال» أو «تحذير»، وفعّل تحذير الزلازل/التحذير الطارئ (Android)؛ أما iPhone فثبّت تطبيق تحذير زلازل موثوقًا واسمح بالإشعارات.
  \item تابع حساب مكتب الأرصاد أو إدارة الطوارئ المحلية على WeChat/الخدمة الرسمية، واسمح بالإشعارات.
  \item تأكد من أن صوت التحذير لن يُكتم تمامًا بوضع عدم الإزعاج؛ وافحص هاتف الوالدين بنفس الطريقة.
\end{enumerate}

\end{tcbpractice}

\booksection{ثامنًا: التعاون عبر الأجهزة: التحكم عن بُعد والنقل السريع وأمان الحسابات}

يكون الهاتف سلسًا غالبًا لأنه «يتعاون» مع الكمبيوتر والجهاز اللوحي. يكمل هذا القسم عدة أدوات عملية عالية الاستخدام: التحكم عن بُعد من الهاتف في الكمبيوتر، ومكتب الذكاء الاصطناعي والتعاون عن بُعد، والخلفيات الديناميكية، والنقل السريع على الشبكة المحلية، والحافظة عبر الأجهزة، ورموز 2FA، وربط المسح الورقي بالمهام والجدول. للتكامل المنهجي عبر الأجهزة راجع الفصل 17.

\booksubsection{التحكم عن بُعد من الهاتف في الكمبيوتر: Sunlogin / ToDesk}

\begin{itemize}
  \item Sunlogin الفئة المستهدفة: من يحتاجون إلى تشغيل Windows أو Mac في المنزل/العمل أثناء الغياب؛ ويُستخدم أيضًا لمساعدة كبار السن عن بُعد (راجع الفصل 27). المزايا الأساسية: منظومة محلية ناضجة، تطبيق الهاتف يتحكم في سطح مكتب الكمبيوتر وينقل الملفات ويشغل/يوقف الجهاز (بعضها يحتاج إلى عتاد)؛ النسخة المجانية تكفي عادةً للاستخدام الشخصي الخفيف.
  \item ToDesk الفئة المستهدفة: من يبحثون عن اتصال مستقر وزمن استجابة منخفض للعمل/الصيانة عن بُعد. المزايا الأساسية: الاتصال برمز الجهاز + كلمة المرور، والتحكم من الهاتف في الكمبيوتر وبالعكس (حسب الإصدار)، ونقل الملفات وتبديل الشاشات المتعددة بسلاسة، وتجربة العقد المحلية عادةً جيدة.
  \item الاستخدام الأساسي: ثبّت العميل على الكمبيوتر وسجّل الدخول/أبقِه متصلًا ثبّت التطبيق على الهاتف أنشئ جلسة برمز الجهاز + كلمة مرور الاتصال (أو نفس الحساب) استخدم لوحة اللمس/الإيماءات للتحكم.
  \item خطوط حمراء أمنية (احفظها جيدًا): إذا ادّعى شخص غريب أنه «خدمة العملاء/الشرطة/النيابة/المحكمة/شركة الاتصالات» وطلب منك تثبيت Sunlogin أو ToDesk والتحكم في هاتفك أو جهاز الكمبيوتر، ارفض دائمًا — هذه طريقة احتيال نموذجية (راجع الفصل 27).
  \item عند الاستخدام الشخصي: فعّل التحقق الثنائي، وغيّر كلمة المرور الافتراضية، وسجّل الخروج عند عدم الاستخدام أو أوقف وضع التشغيل دون مراقبة؛ واتبع قواعد أمان الشركة على أجهزة العمل.
  \item عند العمل عن بُعد في الأماكن العامة، انتبه لمن قد يتجسس على الشاشة وكلمة المرور.
\end{itemize}

\booksubsection{برامج مكتب الذكاء الاصطناعي: Trae / WorkBuddy / Dumate}

\begin{itemize}
  \item Trae الفئة المستهدفة: المستخدمون في البرمجة وتعاون المستندات، ودود للبيئة الصينية. المزايا الأساسية: مساعد برمجة/IDE من ByteDance، وضع SOLO ينشئ مشروعًا بجملة واحدة؛ كفاءة عالية في كتابة الكود وتعديل المستندات. عند التشغيل على الكمبيوتر يمكن متابعة المهام من الهاتف عبر التحكم عن بُعد أو القدرات الرسمية (وفق أحدث أشكال المنتج). راجع الفصل 41 «شمول أدوات البرمجة بالذكاء الاصطناعي».
  \item WorkBuddy الفئة المستهدفة: الموظفون الذين يريدون مساعدة الذكاء الاصطناعي في تلخيص محاضر الاجتماعات والرسائل والتقارير والمستندات المكتبية. المزايا الأساسية: مساعد ذكاء اصطناعي لمواقف العمل، مناسب لـ«المسودة—الصقل—التلخيص»؛ بعض سير العمل يمكن بدؤه من الهاتف، والمهام الثقيلة على الكمبيوتر.
  \item Dumate الفئة المستهدفة: الفرق/الأفراد الذين يحتاجون إلى تعاون متعدد الأشخاص وإدارة موحدة للمستندات والمعرفة. المزايا الأساسية: أداة مكتبية تعاونية بالذكاء الاصطناعي، تركز على تجميع المواد والتحرير التعاوني؛ يمكن العرض والتحرير الخفيف من الهاتف، ويُفضل التعقيد على شاشة الكمبيوتر الكبيرة.
  \item فهم «التحكم عن بُعد من الهاتف في الكمبيوتر»: كثير من قدرات مكتب الذكاء الاصطناعي تعمل على الكمبيوتر/الويب؛ عبر Sunlogin/ToDesk من الهاتف، أو عبر تطبيق/موقع المنتج نفسه، يمكنك «التحكم» في مهام الذكاء الاصطناعي على الكمبيوتر أثناء التنقل. المبدأ: العمليات الثقيلة والمستندات الطويلة على الكمبيوتر؛ ومتابعة التقدم وعدّل الجمل على الهاتف.
  \item تذكير الخصوصية: قبل رفع العقود والهوية وبيانات العملاء، اقرأ سياسة خصوصية المنتج ومتطلبات الامتثال المؤسسي؛ وفضل أدوات العمل المحددة للشركة.
\end{itemize}

\booksubsection{سطح المكتب الاصطناعي: خلفيات miHoYo الديناميكية (هاتف + PC)}

\begin{itemize}
  \item Artificial Desktop الفئة المستهدفة: محبو شخصيات miHoYo ومن يريدون خلفيات ديناميكية مجانية. المزايا الأساسية: تطبيق خلفيات ديناميكية مجاني من miHoYo، يدعم الهاتف وPC؛ يمكن استخدام عروض Live2D/الديناميكية كخلفية، وبعض الطرازات/الأنظمة تدعمها كشاشة قفل أو خلفية سطح مكتب.
  \item توفير الطاقة والأداء: الخلفيات الديناميكية تستهلك بطارية وGPU أكثر من الصور الثابتة؛ يمكن استخدامها فقط أثناء الشحن على الأجهزة القديمة أو المنخفضة الشحن، أو خفض معدل الإطارات (إذا وفر التطبيق الخيار).
  \item مصدر التثبيت: فضّل الموقع الرسمي/المتاجر الرسمية، وتجنب «النسخ المكسورة» التابعة التي ترفق برمجيات خبيثة؛ على PC انتبه لمسار التثبيت وضرورة التشغيل التلقائي.
\end{itemize}

\booksubsection{النقل السريع للملفات على الشبكة المحلية: LocalSend / Feige / AirDrop}

\begin{itemize}
  \item LocalSend الفئة المستهدفة: من يخلطون بين Android/Windows/iOS/Mac/Linux ويريدون نقلًا سريعًا عبر المنصات دون سحابة عامة. المزايا الأساسية: مفتوح المصدر، اتصال مباشر على الشبكة المحلية، نقل الصور والفيديو والمستندات على نفس Wi-Fi، دون كابل ودون خوادم الشركات، والسرعة تعتمد على الموجّه.
  \item Feige ChuanShu (الحمام الزاجل) — لمن؟ للربط بين حواسيب Windows في المكتب وسيناريوهات العمل عبر الشبكة المحلية القديمة. أهم المزايا: أداة عريقة لنقل الملفات والدردشة عبر الشبكة المحلية، مناسبة لنقل الملفات الكبيرة داخل الشبكة الداخلية؛ انتبه إلى تنزيلها من قنوات رسمية، والتزم بسياسات تقنية المعلومات المعتمدة في شبكة الشركة.
  \item AirDrop الفئة المستهدفة: مستخدمو iPhone وiPad وMac. المزايا الأساسية: مشاركة بنقرة واحدة في منظومة Apple، اكتشاف عبر Bluetooth + نقل عبر Wi-Fi مباشر، دون استهلاك بيانات؛ يمكن ضبطه على «جهات الاتصال فقط» لتقليل إزعاج الغرباء.
  \item إضافات أخرى: Mi Share وHuawei Share وvivo Share وغيرها سريعة بين أجهزة العلامة نفسها؛ «مساعد نقل الملفات» في WeChat مريح لكنه محدود بحجم الملف وجهاز الدخول، وللملفات الكبيرة يُفضل LocalSend/AirDrop.
  \item نصائح الاستخدام: اربط كلا الطرفين بنفس Wi-Fi؛ نقل الملفات الحساسة عبر Wi-Fi عام محفوف بالمخاطر، ففضل النقطة الساخنة الخاصة؛ وبعد النقل تأكد من اكتمال الملف عند المستلم.
\end{itemize}

\booksubsection{مزامنة الحافظة عبر الأجهزة: KDE Connect / نقل بين الهاتف والكمبيوتر}

\begin{itemize}
  \item «الحافظة الشاملة» من Apple الفئة المستهدفة: مستخدمو iPhone + Mac. المزايا الأساسية: بعد تسجيل الدخول بنفس Apple ID وتفعيل الاستمرارية، انسخ من الهاتف الصق في Mac (والعكس)، مع النص وبعض الملفات.
  \item KDE Connect الفئة المستهدفة: Android + Windows/Linux/Mac، ومن يريدون مزامنة حافظة مفتوحة المصدر وعكس الإشعارات والإدخال عن بُعد. المزايا الأساسية: بعد الاقتران على الشبكة المحلية تُزامن الحافظة، وتُرسل إشعارات الهاتف إلى الكمبيوتر، ويستخدم الهاتف كمبدل شرائح/لوحة لمس؛ البيانات لا تمر عبر حسابات طرف ثالث.
  \item حلول «نقل بين الهاتف والكمبيوتر» من العلامات التجارية: Huawei Multi-Screen Collaboration/Super Device وXiaomi MiaoXiang وvivo Share وOPPO Connect وHonor Connect وغيرها، غالبًا ما تتضمن حافظة عبر الشاشات وسحب وإفلات الملفات، والأفضل بين الأجهزة من العلامة نفسها.
  \item القناة الرسمية Windows + Android: «اتصال الهاتف / Link to Windows» (تطوير Your Phone) يتيح عرض إشعارات الهاتف والصور على الكمبيوتر، والنسخ واللصق في السيناريوهات المدعومة؛ وبعض الطرازات مثل Samsung تتكامل بعمق أكبر.
  \item تذكير أمني: قد تحتوي الحافظة على كلمات مرور ورموز تحقق؛ استخدم المزامنة العامة بحذر على أجهزة الكمبيوتر المشتركة، وتجنب اللصق مرتين في حقول كلمة المرور؛ وبعض برامج الإدخال تقرأ الحافظة أيضًا، فانتبه للأذونات.
\end{itemize}

\booksubsection{رموز 2FA الديناميكية: Aegis / Microsoft Authenticator}

\begin{itemize}
  \item لماذا تستخدم مصادقة الرموز بدلًا من الرسائل النصية: قد تتعرض رموز التحقق النصية لخطف بطاقة SIM؛ بينما تُولَّد رموز TOTP محليًا حسب خوارزمية الوقت، دون الاعتماد على إشارة الهاتف، وهي أكثر مقاومة للتصيد والاعتراض (مع الحاجة لحماية من البرمجيات الخبيثة ولقطات الشاشة).
  \item Aegis Authenticator الفئة المستهدفة: مستخدمو Android الذين يفضلون المصدر المفتوح والتحكم المحلي. المزايا الأساسية: مجاني ومفتوح المصدر، يدعم النسخ الاحتياطي المشفر والاستيراد والتصدير، وواجهة نظيفة، ويمكن فتح قائمة الرموز بالبصمة/الوجه.
  \item Microsoft Authenticator الفئة المستهدفة: مستخدمو حسابات Microsoft وOffice 365/الحسابات المؤسسية؛ عبر iOS/Android. المزايا الأساسية: الموافقة على تسجيل الدخول بالإشعارات، ومراقبة كلمات المرور والنسخ الاحتياطي للحسابات (حسب إصدار المنتج)، وتكامل وثيق مع Windows/أنظمة الهوية المؤسسية.
  \item Google Authenticator / 2FAS وغيرها: شائعة أيضًا في المنظومة؛ المهم هو وجود نسخ احتياطي مشفر حتى لا تفقد الرموز عند تغيير الجهاز.
  \item إلزامي عند تغيير الجهاز: انقل من الجهاز القديم أو استخدم «رموز احتياطية/مفاتيح استعادة» من الخدمة؛ اطبع الرموز الاحتياطية على ورق واحفظها في مكانين، ولا تكتفِ بجهاز واحد.
  \item مع مدير كلمات المرور: بعض خزائن كلمات المرور تخزن TOTP أيضًا، وهو مريح لكن المخاطر مركزة؛ للحسابات عالية القيمة يُفضل «مدير كلمات مرور + مصادق مستقل». راجع الفصل 14.
\end{itemize}

\booksubsection{التركيبة الموصى بها: كيف تربط الأدوات في يوم واحد}

\textbf{مثال على سير العمل العملي}\par

\begin{itemize}
  \item الصباح: راجع جدول اليوم في TickTick اطلّع على الطقس/التحذيرات افتح Trae/المستندات على الكمبيوتر، وضع الهاتف جانبًا.
  \item النهار: حوّل المواد الورقية إلى PDF عبر تطبيقات المسح، وأرسلها إلى الكمبيوتر عبر LocalSend/AirDrop؛ اكتب المسودات الطويلة على الكمبيوتر بأدوات مكتب الذكاء الاصطناعي.
  \item أثناء الخروج: اتصل بالكمبيوتر عبر ToDesk/Sunlogin لمتابعة المهام؛ وسجّل الدخول إلى الحسابات المهمة عبر رموز Aegis/Microsoft Authenticator.
  \item المساء: علّق المهام المكتملة في TickTick؛ وإذا احتجت غيّر الخلفية عبر Artificial Desktop للاسترخاء.
\end{itemize}

\begin{tcbwarn}{قواعد التحكم عن بُعد وأمان الحسابات}
\begin{itemize}
  \item إذا طلب منك شخص تثبيت برنامج تحكم عن بُعد أغلق الهاتف، واتصل بخدمة العملاء الرسمية بنفسك.
  \item لا تخبر أحدًا برموز التحقق أو الرموز الديناميكية، حتى «التوصيل» أو «خدمة العملاء».
  \item قبل التحكم عن بُعد في أجهزة العمل أو تثبيت برامج الذكاء الاصطناعي، اسأل قسم تكنولوجيا المعلومات أولًا.
  \item كن حذرًا جدًا عند نقل الملفات أو تسجيل الدخول عن بُعد عبر الشبكات العامة.
\end{itemize}

\end{tcbwarn}

\begin{tcbpractice}{تطبيق عملي: شغّل رابط «الأجهزة المتعددة» مرة واحدة}
\begin{enumerate}
  \item اربط الهاتف والكمبيوتر بنفس Wi-Fi، وانقل صورة عبر LocalSend (أو AirDrop).
  \item أنشئ مهمة في TickTick: «الساعة 8 مساءً: ترتيب المسح الضوئي»، وتأكد من ظهورها على الهاتف والكمبيوتر.
  \item فعّل 2FA لبريدك الإلكتروني أو WeChat (إذا كان مدعومًا)، واربطه بـ Aegis أو Microsoft Authenticator، واحفظ الرموز الاحتياطية.
  \item ثبّت Sunlogin أو ToDesk على جهاز الكمبيوتر الخاص بك، واتصل بنجاح من هاتفك مرة واحدة، ثم اخرج أو قوِّ كلمة المرور، لتعرف العملية دون ترك ثغرة أمنية.
\end{enumerate}

\end{tcbpractice}
